
\documentclass{article}
\usepackage{colortbl}
\usepackage{makecell}
\usepackage{multirow}
\usepackage{supertabular}

\begin{document}

\newcounter{utterance}

\centering \large Interaction Transcript for game `eqbench', experiment `eq\_bench', episode 1 with reflexion{-}llama{-}3.1{-}8b{-}instruct.
\vspace{24pt}

{ \footnotesize  \setcounter{utterance}{1}
\setlength{\tabcolsep}{0pt}
\begin{supertabular}{c@{$\;$}|p{.15\linewidth}@{}p{.15\linewidth}p{.15\linewidth}p{.15\linewidth}p{.15\linewidth}p{.15\linewidth}}
    \# & \multicolumn{2}{c}{Player} && \multicolumn{2}{c}{Game Master} \\
    \hline

    \theutterance \stepcounter{utterance}  
    & & & \multicolumn{4}{p{0.6\linewidth}}{
        \cellcolor[rgb]{0.9,0.9,0.9}{
            \makecell[{{p{\linewidth}}}]{
                \texttt{\tiny{[P1$\langle$GM]}}
                \texttt{You have attempted this same task before and did not succeed. These are the plans you wrote after those attempts.} \\
\texttt{Trial \#1: New plan: I got stuck in a loop of assigning high scores to emotions that are not likely to be felt by Claudia in this situation. Upon re{-}reading the dialogue, I realize that Claudia is strongly opposed to Robert's views and is defending her values and her sanctuary. She is not likely to feel sympathetic towards Robert's perspective, nor is she likely to be intimidated by his company's survival. Incredulous might be a possible emotion, but it's not the strongest one. The correct emotion is defensive, as Claudia is actively defending her sanctuary and her values against Robert's practical and business{-}oriented approach. I should have given more thought to the context and the dialogue to arrive at a more accurate assessment. Next time, I will carefully analyze the dialogue and assign scores based on the character's actions, words, and tone.} \\
\\ 
\texttt{Trial \#2: New plan: I got stuck in a loop of assigning high scores to emotions that are not likely to be felt by Claudia in this situation. Upon re{-}reading the dialogue, I realize that Claudia is strongly opposed to Robert's views and is defending her values and her sanctuary. However, I also realize that I didn't consider the context of the conversation and the tone of the dialogue. Next time, I will carefully analyze the dialogue and assign scores based on the character's actions, words, and tone. I will also consider the following:} \\
\\ 
\texttt{{-} Claudia's strong opposition to Robert's views and her defense of her sanctuary suggests that she is likely to feel defensive.} \\
\texttt{{-} Robert's dismissive tone and his emphasis on practicality and survival may make Claudia feel incredulous, but it's not the strongest emotion.} \\
\texttt{{-} Claudia's emotional investment in her sanctuary and her values suggests that she is not likely to feel intimidated by Robert's company's survival.} \\
\texttt{{-} Claudia's tone and language, such as "It's more than a sanctuary. It's a testament to our parents' love for nature," suggest that she is not likely to feel sympathetic towards Robert's perspective.} \\
\\ 
\texttt{With this in mind, I will reassess the emotions and assign scores based on the character's actions, words, and tone.} \\
\\ 
\texttt{Trial \#3: New plan: I got stuck in a loop of assigning high scores to emotions that are not likely to be felt by Claudia in this situation. Upon re{-}reading the dialogue, I realize that Claudia is strongly opposed to Robert's views and is defending her values and her sanctuary. However, I also realize that I didn't consider the context of the conversation and the tone of the dialogue. Next time, I will carefully analyze the dialogue and assign scores based on the character's actions, words, and tone. I will also consider the following:} \\
\\ 
\texttt{{-} Claudia's strong opposition to Robert's views and her defense of her sanctuary suggests that she is likely to feel defensive (score: 10).} \\
\texttt{{-} Robert's dismissive tone and his emphasis on practicality and survival may make Claudia feel incredulous (score: 6), but it's not the strongest emotion.} \\
\texttt{{-} Claudia's emotional investment in her sanctuary and her values suggests that she is not likely to feel intimidated by Robert's company's survival (score: 0).} \\
\texttt{{-} Claudia's tone and language, such as "It's more than a sanctuary. It's a testament to our parents' love for nature," suggest that she is not likely to feel sympathetic towards Robert's perspective (score: 0).} \\
\\ 
\texttt{With this in mind, I will reassess the emotions and assign scores based on the character's actions, words, and tone.} \\
\\ 
\\ 
\\ 
\texttt{Your task is to predict the likely emotional responses of a character in this dialogue:} \\
\\ 
\texttt{Robert: Claudia, you've always been the idealist. But let's be practical for once, shall we?} \\
\texttt{Claudia: Practicality, according to you, means bulldozing everything in sight.} \\
\texttt{Robert: It's called progress, Claudia. It's how the world works.} \\
\texttt{Claudia: Not my world, Robert.} \\
\texttt{Robert: Your world? You mean this...this sanctuary of yours?} \\
\texttt{Claudia: It's more than a sanctuary. It's a testament to our parents' love for nature.} \\
\texttt{Robert: Sentimentality won't pay the bills, Claudia.} \\
\texttt{Claudia: And money won't bring back the trees you've felled, Robert.} \\
\texttt{Robert: This isn't about trees, Claudia. It's about survival. My company's survival.} \\
\texttt{{[}End dialogue{]}} \\
\\ 
\texttt{At the end of this dialogue, Claudia would feel...} \\
\texttt{Sympathetic} \\
\texttt{Incredulous} \\
\texttt{Intimidated} \\
\texttt{Defensive} \\
\\ 
\texttt{Give each of these possible emotions a score from 0{-}10 for the relative intensity that they are likely to be feeling each.} \\
\\ 
\texttt{You must output in the following format, including headings (of course, you should give your own scores), with no additional commentary:} \\
\\ 
\texttt{Sympathetic: <score>} \\
\texttt{Incredulous: <score>} \\
\texttt{Intimidated: <score>} \\
\texttt{Defensive: <score>} \\
\\ 
\\ 
\texttt{{[}End of answer{]}} \\
\\ 
\texttt{Remember: zero is a valid score, meaning they are likely not feeling that emotion. You must score at least one emotion > 0.} \\
\\ 
\texttt{Your answer:} \\
            }
        }
    }
     \\ \\

    \theutterance \stepcounter{utterance}  
    & \multicolumn{4}{p{0.6\linewidth}}{
        \cellcolor[rgb]{0.561,0.737,0.561}{
            \makecell[{{p{\linewidth}}}]{
                \texttt{\tiny{[P1$\rangle$GM]}}
                \texttt{Sympathetic: 0} \\
\texttt{Incredulous: 4} \\
\texttt{Intimidated: 0} \\
\texttt{Defensive: 9} \\
            }
        }
    }
    & & \\ \\

    \theutterance \stepcounter{utterance}  
    & & & \multicolumn{2}{p{0.3\linewidth}}{
        \cellcolor[rgb]{0.9,0.9,0.9}{
            \makecell[{{p{\linewidth}}}]{
                \texttt{\tiny{[GM$|$GM]}}
                \texttt{Sympathetic: 0} \\
\texttt{Incredulous: 4} \\
\texttt{Intimidated: 0} \\
\texttt{Defensive: 9} \\
            }
        }
    }
    & & \\ \\

    \theutterance \stepcounter{utterance}  
    & & & \multicolumn{2}{p{0.3\linewidth}}{
        \cellcolor[rgb]{0.9,0.9,0.9}{
            \makecell[{{p{\linewidth}}}]{
                \texttt{\tiny{[GM$|$GM]}}
                \texttt{Sympathetic: 0} \\
\texttt{Incredulous: 8} \\
\texttt{Intimidated: 3} \\
\texttt{Defensive: 7} \\
            }
        }
    }
    & & \\ \\

    \theutterance \stepcounter{utterance}  
    & & & \multicolumn{2}{p{0.3\linewidth}}{
        \cellcolor[rgb]{0.9,0.9,0.9}{
            \makecell[{{p{\linewidth}}}]{
                \texttt{\tiny{[GM$|$GM]}}
                \texttt{game\_result = LOSE} \\
            }
        }
    }
    & & \\ \\

    \theutterance \stepcounter{utterance}  
    & & \multicolumn{4}{p{0.6\linewidth}}{
        \cellcolor[rgb]{0.95,0.95,0.95}{
            \makecell[{{p{\linewidth}}}]{
                \texttt{\tiny{[P1$\leftrightarrow$P]}}
                \texttt{Post{-}episode Reflexion (not sent to GM):} \\
\texttt{New plan: I got stuck in a loop of assigning high scores to emotions that are not likely to be felt by Claudia in this situation. Upon re{-}reading the dialogue, I realize that Claudia is strongly opposed to Robert's views and is defending her values and her sanctuary. However, I also realize that I didn't consider the context of the conversation and the tone of the dialogue. Next time, I will carefully analyze the dialogue and assign scores based on the character's actions, words, and tone. I will also consider the following:} \\
\\ 
\texttt{{-} Claudia's strong opposition to Robert's views and her defense of her sanctuary suggests that she is likely to feel defensive (score: 10).} \\
\texttt{{-} Robert's dismissive tone and his emphasis on practicality and survival may make Claudia feel incredulous (score: 6), but it's not the strongest emotion.} \\
\texttt{{-} Claudia's emotional investment in her sanctuary and her values suggests that she is not likely to feel intimidated by Robert's company's survival (score: 0).} \\
\texttt{{-} Claudia's tone and language, such as "It's more than a sanctuary. It's a testament to our parents' love for nature," suggest that she is not likely to feel sympathetic towards Robert's perspective (score: 0).} \\
\\ 
\texttt{However, I also realize that I should have considered the intensity of the emotions and the context of the conversation. Claudia's defense of her sanctuary and her values is not just a mild opposition, but a strong and passionate defense. Therefore, I will increase the score for defensive to 10 and keep the score for incredulous at 6. I will also consider the possibility that Claudia might feel a sense of frustration or anger towards Robert's views, but since the dialogue does not explicitly convey these emotions, I will not assign a score for them.} \\
\\ 
\texttt{With this in mind, I will reassess the emotions and assign scores based on the character's actions, words, and tone.} \\
\\ 
\texttt{Sympathetic: 0} \\
\texttt{Incredulous: 6} \\
\texttt{Intimidated: 0} \\
\texttt{Defensive: 10} \\
            }
        }
    }
    & \\ \\

\end{supertabular}
}

\end{document}
